Our work builds on a large body of research on CGH algorithms, which we review below. For a more comprehensive overview of holographic displays, we refer the reader to~\cite{yaracs2010state,park2017recent,chang2020toward,javidi2021roadmap, pi2022review}.

\paragraph{\bf{CGH Algorithms}}
Holograms create a visible image or 3D scene indirectly by displaying an interference pattern, i.e., the hologram, on a 2D amplitude- or phase-only SLM. Methods that convert a target intensity distribution into a hologram are called CGH algorithms. These algorithms have been developed to accommodate a wide variety of input 3D representations, including point clouds~\cite{h2009computer, lucente1993interactive}, meshes~\cite{matsushima2003fast, ahrenberg2008computer}, wireframes~\cite{blinder2021real}, light fields~\cite{padmanaban2019holographic, zhang2019three, park2019non, kang2008accurate, choi2022time}, image layers~\cite{chen2015improved, shi2022end}, and most recently, Gaussians~\cite{choi2025gaussian}. Classic and deep learning–based \textit{direct} CGH methods follow a common pipeline: the target intensity is first encoded into a 2D complex wavefront, typically assuming a certain phase distribution; this wavefront is then propagated to the SLM plane via models such as the angular spectrum method~\cite{goodman2005introduction, pellat1994fresnel}; finally, the resulting complex field is converted into a phase- or amplitude-only pattern, depending on the SLM type~\cite{tsang2013novel, maimone2017holographic, qi2016speckleless}. In contrast, \emph{iterative} CGH methods make use of iterative optimization to achieve a better image quality, albeit at the cost of increased runtime~\cite{gerchberg1972practical, fienup1980iterative, zhang20173d, peng2020neural, chakravarthula2019wirtinger}.

\paragraph{\bf{Phase Distributions of Holograms}}

Although phase is not directly observable, the phase profile of a wavefront plays a crucial role in determining the spatio-angular behavior of the observable light field~\cite{st1995phase, schiffers2023stochastic, kim2022accommodative, Chakravarthula2022pupil}. For this reason, two popular heuristics have been developed that are widely used in  CGH literature: \textit{smooth-phase} and \textit{random-phase} holograms~\cite{maimone2017holographic, yoo2021optimization}. Smooth-phase, sometimes called random-phase-free, holograms~\cite{shimobaba2015random} achieve high image quality that can be demonstrated with relatively simple experimental setups~\cite{maimone2017holographic, shi2021towards, peng2020neural, choi2025gaussian}. The main drawback of smooth-phase distributions, however, is that they concentrate energy in the low frequencies of the angular spectrum, resulting in a severely restricted eye box size, limited defocus effects, and increased sensitivity to pupil position. These effects limit the perceptual realism and overall user experience of the produced holograms~\cite{kim2024holographic} as well as the support for perceptually important focus cues~\cite{kim2022accommodative}.

On the other hand, random-phase holograms are capable of reconstructing larger parallax and natural defocus blur, which is necessary for a perceptually realistic and natural viewing experience \cite{kim2024holographic}. However, rapid phase variations between adjacent pixels introduce unwanted speckle noise created by constructive and destructive interference~\cite{goodman2007speckle}. Speckle reduction techniques often utilize some form of partial coherence, introduced by partially coherent or multiple coherent light sources, or more commonly through time multiplexing~\cite{lee2020light, chao2024large, kuo2023multisource, choi2022time, curtis2021dcgh, peng2021speckle}. 

In this paper, we prove that only a specific family of random-phase distributions yields correct alpha blending of wavefronts (Sec. \ref{sec:awb_random}). We then show that the conventional uniform $2\pi$ distribution used in prior random-phase literature conveniently happens to be one of them, allowing our RPWS holograms to maximally utilize the SLM bandwidth for large parallax and natural defocus reconstruction. Finally, we show that time multiplexing is not only a speckle reduction heuristic for random-phase holography, but the statistically exact mechanism that enables accurate alpha blending.



\paragraph{\bf{Primitive-based 3D Representations and CGH Algorithms.}} Among the various CGH algorithms described above, polygon-based CGH algorithms that use meshes as the input 3D format are closely related to our approach \cite{matsushima2005computer, matsushima2005exact, matsushima2009extremely, askari2017occlusion, chen2009computer, matsushima2018full, matsushima2014silhouette}. However, the constraint that each triangle has to be fully opaque limits the expressivity of mesh models by preventing soft blending of primitives, a paradigm central to recent 3D reconstruction frameworks such as 3DGS \cite{kerbl3Dgaussians}. Coupled with per-primitive texturing, these restrictions make polygon-based CGH difficult to capture fine details without using an excessive number of tiny triangles. Gaussian Wave Splatting \cite{choi2025gaussian}, on the other hand, has emerged as a new primitive-based CGH algorithm that converts Gaussian-based scene representations \cite{kerbl3Dgaussians, huang20242d} to complex holograms, greatly outperforming traditional polygon-based CGH in terms of image quality and rendering time. However, GWS does not support view-dependent effects and natural defocus due to its smooth-phase nature. Additionally, the alpha-blending and wave-splatting procedure described in GWS does not naturally extend to random-phase Gaussians. Na\"ively applying GWS to random-phase wavefronts, as suggested in the GWS supplementary materials (which we refer to as GWS-NRP), provably fails as we demonstrate throughout the paper. Finally, GWS solely works with Gaussians, thus it remains unknown if such a wave compositing scheme can be extended to support a larger variety of translucent geometric primitives widely used in recent 3D reconstruction literature, such as soft-edged triangles \cite{Held2025Triangle, sheng20252dtrianglesplattingdirect, burgdorfer2025radianttrianglesoupsoft}.

In this work, we show that for random-phase primitive wavefronts, alpha blending is exact \textit{in expectation} in the \textit{intensity domain}, in contrast to the amplitude domain for smooth-phase primitives as described in GWS (Sec. \ref{sec:awb_random}). We then devise a novel wavefront composition technique that specifically works with arbitrary random-phase geometric primitives (Sec. \ref{subsec:gws_rp}) for accurate color reconstruction and occlusion handling. As such, our algorithm generates holograms with high image quality by fully leveraging the translucent primitive representations while synthesizing accurate defocus blur and wide parallax via random phase and time-multiplexing.
